\documentclass[tikz,border=6pt]{standalone}
\usepackage{ctex}
\usepackage{amsmath}
\usepackage{amssymb}
\usepackage{pgfplots}
\pgfplotsset{compat=1.18}
\usetikzlibrary{calc,angles,quotes,arrows.meta,positioning,decorations.markings}
\begin{document}
\begin{tikzpicture}
\begin{axis}[
    width=11cm, height=9cm,
    xlabel={$x$}, ylabel={$y=\log_a x$},
    xmin=-0.15, xmax=8.3, ymin=-3.4, ymax=3.4,
    axis lines=middle,
    xtick={1,2,3,4,5,6,7,8},
    ytick={-3,-2,-1,1,2,3},
    every axis x label/.style={at={(ticklabel cs:1.0)},anchor=west},
    every axis y label/.style={at={(ticklabel cs:1.0)},anchor=south},
    tick label style={font=\footnotesize},
    label style={font=\small},
    samples=300,
    clip=false,
]
% a=2 (递增, 最陡)
\addplot[blue,very thick,domain=0.05:8] {ln(x)/ln(2)};
% a=e (递增)
\addplot[teal,very thick,domain=0.05:8] {ln(x)};
% a=10 (递增, 最缓)
\addplot[orange!80!black,very thick,domain=0.05:8] {ln(x)/ln(10)};
% a=1/2 (递减)
\addplot[red,very thick,domain=0.05:8] {ln(x)/ln(0.5)};

% 公共点 (1,0)
\addplot[only marks,mark=*,black,mark size=2pt] coordinates {(1,0)};
\node[anchor=north west,font=\scriptsize] at (axis cs:1.05,-0.08) {$(1,0)$};

% 曲线标签
\node[blue,font=\small] at (axis cs:5.0,2.55) {$a=2$};
\node[teal,font=\small] at (axis cs:7.2,2.15) {$a=e$};
\node[orange!80!black,font=\small] at (axis cs:7.4,1.05) {$a=10$};
\node[red,font=\small] at (axis cs:5.6,-2.55) {$a=\tfrac12$};

% 注释框
\node[fill=white,fill opacity=0.85,text opacity=1,draw=blue!40,rounded corners,
      font=\scriptsize,align=left,inner sep=4pt] at (axis cs:2.55,-2.35)
  {$a>1$：递增曲线\\（底越小越陡）};
\node[fill=white,fill opacity=0.85,text opacity=1,draw=red!40,rounded corners,
      font=\scriptsize,align=left,inner sep=4pt] at (axis cs:2.55,2.55)
  {$0<a<1$：递减曲线};

% 全部过 (1,0) 的强调
\node[fill=white,fill opacity=0.85,text opacity=1,draw=black!50,rounded corners,
      font=\scriptsize,align=center,inner sep=3pt] at (axis cs:6.4,-0.65)
  {所有曲线都\\经过 $(1,0)$};

\end{axis}
\end{tikzpicture}
\end{document}
